
Pour vérifier que la suite est exacte, on peut le faire localement sur @@M0@@, et utiliser les sections définies dans l'énoncé de 2.5. On a alors une base de @@M1@@ formée des images des @@M2@@, une base de @@M3@@ formée des images des @@M4@@, et une base de @@M5@@ formée des images des @@M6@@. Il est alors clair que @@M7@@ est surjectif, et que son noyau est formé des éléments ayant des coordonnées nulles sur les @@M8@@ et les @@M9@@, et des coordonnées égales pour tout @@M10@@ sur @@M11@@ et @@M12@@. Par suite, la suite est également exacte au milieu. Or les trois @@M13@@-Modules intervenant sont localement libres de type fini; comme @@M14@@ est de rang @@M15@@ d'après la démonstration du i), on voit immédiatement que la suite est exacte.

\subsection*{2.6.}
On suppose maintenant que @@M16@@ est un schéma ayant un faisceau ample. Tout sous-schéma fermé @@M17@@ de @@M18@@ a alors un faisceau ample. On sait que dans ce cas ``il existe des résolutions globales'', et que les @@M19@@ formés respectivement avec les faisceaux localement libres de type fini et les complexes parfaits sont identiques (II, 2.2.9). Si en outre l'immersion @@M20@@ est régulière, c'est un morphisme parfait, et on peut définir un homomorphisme additif @@M21@@ (IV 2.12), en associant à la classe d'un complexe parfait @@M22@@ la classe du complexe @@M23@@ (@@M24@@ étant un foncteur exact).

\textbf{Proposition 2.7.} \uline{Soient} @@M25@@ \uline{un schéma ayant un faisceau ample}, @@M26@@ \uline{une immersion fermée régulière}, @@M27@@ \uline{le faisceau conormal de} @@M28@@ \uline{dans} @@M29@@. \uline{Alors}, \uline{pour tout} @@M30@@, \uline{on a la relation}:
@@M31@@
\uline{où} @@M32@@ \uline{est la classe de} @@M33@@ \uline{dans} @@M34@@, \uline{et} @@M35@@.


